\section{The New Dawn}

\begin{quotex}
To defend an ideal and hold positions even if they be lost positions, or better said: even when it may be problematic that those who still remain awake during the night could meet those who will appear in the new dawn.
\flright{\textsc{Julius Evola}}
\end{quotex}


We have been demonstrating the coherence of the Western Tradition and the inner connections starting from the supposed Hyperborean primordial tradition and subsequently to the Vedic, ancient Greco-Roman, and Medieval civilizations. If it is the case, as Nietzsche claims, that everything is a perspective, we can point out that most perspectives are based on rhetoric, but very few on conviction. The former are chatty, on the assumption that the will to power is a battle of words or ideas; hence, they engage in sterile debates, vilify ideological foes, and create shibboleths to identify friends. That is, they look at Tradition from the outside, as just another influence, to be juxtaposed with other seemingly similar thinkers. They gather opinions like squirrels gathering acorns, with little concern for the spirit that gave rise to them.

To be convinced, on the other hand, is to know Tradition from the inside, as the direct intuition of interior states, and to know oneself as the creator as his life. As \textbf{Carlo Michelstaedter} writes of the man of conviction:

\begin{quotex}
He must take on himself the responsibility for his life (as it must be lived in order for him to attain life), which cannot rest with another. He must have in himself the certainty of his own life, which others cannot give him. He must create himself and the world, which does not exist before him: he must be master and not slave in his house.
\end{quotex}

\textbf{Rene Guenon} has pointed out that time is not a sequence of indifferent moments of the same quantity; time is also qualitative. Julius Evola draws the logical conclusion that at a given period of historico-material conditions, different spirits will be attracted to that world; this he calls the \textbf{Law of Elective Affinity}. Hence, there are many spirits who are blind to the teachings of Tradition and experience the modern world as completely rational, good, and progressive. This reflects the formlessness and increasing lack of differentiation of the end of a cycle.

At the same time, the owl of Minerva flies at dusk\footnote{See Section \ref{sec:OwlMinerva} in this book.}. That means that only at the end of the cycle can the whole be seen in retrospect. Men are no longer tied to the Tradition of their land, but are exposed to those of the entire world. They are no longer forced into the role of an apologist for one or another, but can instead focus on the heart of their doctrines. Thus, by the same law mentioned above, men of spirit, who manage to stay vigilant in the dark, will also be attracted to this world.

Some of them may constitute the new elite, if not now, then in the next generation. From where they come, it is not clear. The English speaking world is moribund and is burdened with the absurd notion of the ``marketplace of ideas''. Europe believes in its own decline, but no longer believes in resurrection. It is not easy to predict from where the elite will arise; perhaps a few Magi will be able to read the skies.

There is a German saying: ``Wer sagt A muss B sagen'' (Who says A must say B). We think that through these many translations and commentaries we have said A. Now we want to say B and will focus on some specific consequences. Beyond that, it is unclear. I am awaiting the arrival of the more mystical writings of \textbf{Guido de Giorgio}; if they are of interest, I will communicate them.

The other half of perspective is the will to power; the idea alone is just a possibility. A possibility of manifestation, though, requires power to bring essence into being, against forces of resistance and the inertia and chaos of prime matter. That is entirely different from writing and requires men of sufficient strength and vigor. Perhaps the new dawn has arrived already.


\flrightit{Posted on 2012-04-22 by Cologero}

